\documentclass[10pt,a4paper]{amsart}
\usepackage[margin=19mm]{geometry}
\usepackage{amsmath,amssymb,mathtools}
\usepackage[scr=rsfs,cal=euler]{mathalfa}
\usepackage{xfrac}
\usepackage[unicode,hidelinks,bookmarks=false]{hyperref}
\newcommand{\ie}{i.e.\ }
\newcommand{\cf}{cf.\ }
\newcommand{\resp}{respectively\ }
\newcommand{\Subsection}{\subsection}
\providecommand{\nobreakdash}{\nobreak\hbox{-}}
\setlength{\emergencystretch}{2em}
\multlinegap0pt
\usepackage{amssymb}
\usepackage{mathtools}
\usepackage{xcolor}
\usepackage{enumerate}
\newtheorem{lemma}{Lemma}
\title{Dynamical Mordell--Lang conjecture for split self-maps of affine curve times projective curve}
\author{Junyi Xie, She Yang, Aoyang Zheng}
\date{Source: arXiv:2602.08608v2 (CC BY 4.0)}
\begin{document}
\maketitle

\noindent\emph{Source and licence.} Dynamical Mordell--Lang conjecture for split self-maps of affine curve times projective curve. Version arXiv:2602.08608v2 (CC BY 4.0), distributed by arXiv under the Creative Commons Attribution 4.0 International licence, http://creativecommons.org/licenses/by/4.0/, as recorded in the arXiv licence metadata for this exact version and shown on the abstract page https://arxiv.org/abs/2602.08608v2. Full licence notice: \texttt{licenses/arxiv-2602.08608-CC-BY-4.0.txt}.\par\smallskip

\noindent\textit{Editorial scope.} Counted unit: the article's lemma growth (source0.tex lines 195--198). The article's complete original proof of it is delivered with it (source0.tex lines 200--213, one \texttt{proof} environment of 14 lines). No mathematics is written, repaired or re-derived: the statement and the proof are the article's own lines, in the article's own notation, and no retained line is substituted or re-flowed.  No further source-local context is needed: every cross-reference of every retained line resolves inside the two delivered spans.  Nothing was omitted from any retained span, so the package carries no omission record.  No retained line contains a citation, so the wrapper carries no bibliography and no bibliography entry.  No retained line contains a figure environment, an image-inclusion command or drawing code, so no image is delivered and the package declares no asset dependency.  Editorial formatting only, in addition: an AMS \texttt{amsart} preamble that restates the article's own declarations and macros used by the retained lines, namely the environments \texttt{lemma} on separate counters, together with the standard packages those lines need.  The retained environments print in this document as Lemma 1 in order of appearance, whereas the article prints the same items with the numbering of its own full document.  The complete native source of the article as distributed in the arXiv e-print for this exact version is bundled unchanged as \texttt{sources/194258/source0.tex}.
\begin{lemma}\label{norm growth}
    If $x_0 \in \mathbb{A}^1(K)$ is not a preperiodic point of f, then there exist a place $v \in M_K$, constants $c_1,c_2>0$, and a positive integer $N$, such that for $n>N$, we have
    $$ c_1 d^n < \log |f^n(x_0)|_v < c_2 d^n. $$
\end{lemma}

\begin{proof}
    Write $f = a_d x^d + a_{d-1} x^{d-1} + \cdots + a_0$, where $a_d \neq 0$. Denote $M_{K,\infty}$ as the set of archimedean places of $K$. Let $S = M_{K,\infty}\cup\{ v\in M_K \big{|}\ |a_d|_v \neq 1\}\cup\bigcup\limits_{i=0}^{d-1}\{ v\in M_K \big{|}\ |a_i|_v > 1\}$. Note that $S$ is a finite set. For $v\in S$, we denote $C_v=\frac{2}{|a_d|_v}(1+\sum\limits_{i=0}^{d-1}|a_i|_v)+1$.
    
    Since $O_f(x_0)$ is infinite, we know $\{h(f^n(x_0)) |\ n\in \mathbb{N}\}$ is unbounded by the Northcott property. Here $h$ is the height function. Then we can find either a place $v\notin S$ together with an integer $N$ such that $|f^N(x_0)|_v > 1$, or a place $v\in S$ together with an integer $N$ such that $|f^N(x_0)|_v > C_v$.
    
    In the previous case, we have $|f^{n+1}(x_0)|_v = |f^n(x_0)|_v^d$ when $n\geq N$. Hence the lemma follows.
    
    In the latter case, the inequalities $\frac{1}{2} |a_d|_v < \frac{|f^{n+1}(x_0)|_v}{|f^n(x_0)|_v^d} < \frac{3}{2} |a_d|_v$ and $|f^n(x_0)|_v>C_v$ hold for every $n\geq N$. Taking logarithm, we get
    $$ \log(\frac{1}{2}|a_d|_v) < \log |f^{n+1}(x_0)|_v - d\log|f^n(x_0)| < \log(\frac{3}{2}|a_d|_v).$$
    
    Then for $n\geq N$, by taking summation in the standard way, we get $$(\log|f^N(x_0)|_v+\frac{\log(|a_d|_v/2)}{d-1})d^{n-N} - \frac{\log(|a_d|_v/2)}{d-1} < \log |f^n(x_0)|_v$$
    $$< (\log|f^N(x_0)|_v+\frac{\log(3|a_d|_v/2)}{d-1})d^{n-N} - \frac{\log(3|a_d|_v/2)}{d-1}.$$
    Thus we finish the proof.
\end{proof}

\end{document}
